% ─────────────────────────────────────────────────────────────────────────────
\section{Graybox Transfer Protocol}
\label{sec:graybox-design}
% ─────────────────────────────────────────────────────────────────────────────

To evaluate whether the proposed method changes the structure of learned robustness
rather than only improving whitebox accuracy, we conduct a cross-model graybox transfer
experiment. The experiment uses four method families: PGD-AT,
DDN-AT, Multi-AT, and AttackDRO++ (Ours).
Each method is evaluated over five representative seeds, giving 20 checkpoints in total.

Each checkpoint is used in two roles. As a \emph{surrogate}, it generates adversarial
examples. As a \emph{target}, it is evaluated on adversarial examples generated by all
surrogate checkpoints. Therefore, the full protocol evaluates
\[
    20 \times 20 = 400
\]
surrogate--target pairs. This design produces whitebox, same-method transfer, and
cross-method graybox transfer results from a single unified evaluation table.

To keep the computation tractable while still supporting per-class and per-attack
analysis, each surrogate checkpoint generates a balanced adversarial cache. The CIFAR-10
test set contains 1000 images per class. For each class, the test images are split evenly
across the eight evaluation attacks. Thus, each attack-class cell contains 125 source
images:
\[
    10 \text{ classes}
    \times
    8 \text{ attacks}
    \times
    125 \text{ examples}
    =
    10{,}000
\]
adversarial examples per surrogate checkpoint.

The class-wise split is generated deterministically using a fixed shuffle seed. Therefore,
the same attack-class cell always refers to the same underlying CIFAR-10 source images
across surrogate checkpoints. This makes paired comparisons between surrogate and
target methods well aligned.

For each surrogate--target pair, the evaluation produces one row per attack-class group,
indexed by surrogate method, surrogate seed, target method, target seed, attack, and
class. Since there are 80 attack-class groups per pair, the full table contains
\(400 \times 80 = 32{,}000\) rows. Transfer matrices, attack-class drilldowns, and
paired method comparisons are all computed by aggregating this same long-form table. Table~\ref{tab:graybox_transfer_regimes} summarizes the transfer regimes used in the
analysis.

\begin{table}[H]
\centering
\caption{Transfer regimes in the graybox evaluation protocol.}
\label{tab:graybox_transfer_regimes}
\footnotesize
\setlength{\tabcolsep}{4pt}
\renewcommand{\arraystretch}{1.08}
\begin{tabularx}{\textwidth}{@{}
>{\raggedright\arraybackslash}p{0.25\textwidth}
>{\raggedright\arraybackslash}p{0.25\textwidth}
>{\raggedright\arraybackslash}p{0.17\textwidth}
>{\raggedright\arraybackslash}X
@{}}
\toprule
Regime & Condition & Rows & Purpose \\
\midrule

Whitebox
& Same method and same seed
& $20 \times 80 = 1{,}600$
& Recovers standard whitebox robustness as a sanity check. \\

Same-method transfer
& Same method, different seed
& $80 \times 80 = 6{,}400$
& Measures seed-to-seed transfer within a method family. \\

Cross-method same-seed
& Different methods, same seed
& $60 \times 80 = 4{,}800$
& Measures transfer between method families under matched seeds. \\

Cross-method cross-seed
& Different methods, different seeds
& $240 \times 80 = 19{,}200$
& Measures the most general cross-method graybox transfer setting. \\

\bottomrule
\end{tabularx}
\end{table}

In Chapter~6, the cross-method same-seed and cross-method cross-seed regimes are
combined to form the main graybox transfer analysis. The whitebox and same-method
transfer regimes are reported separately as diagnostic checks.
